% ---------------------------------------------------------------------
% Front matter: abstract (English) and reader's guide (Chinese)
% ---------------------------------------------------------------------

\begin{abstract}
\noindent
The \AEG{} manuscripts read an arithmetic expression through the ordered
histories that evaluate it: an arithmetic expression space carries an
assignment function and a metric, the four operations act as motions, and the
order defect of two histories with equal total charges is an exact weighted
area in the accumulative commutative space (\ACS).  This note asks a single
question of that machinery: what happens if the ideal point $q$ of an
arithmetic history is read as an asymptotic state, the ordinary point $p$ as a
finite-scale dressed state, and the additive and multiplicative charge steps as
the emission and absorption of a gauge quantum?  The answer is developed
section by section, and each step is labelled as a proved \AEG{} statement, a
declared dictionary hypothesis, or a computation carried out here.

Five results are obtained inside \AEG{} and need no dictionary at all.  First,
the family of $B_\infty$-histories joining an ultraviolet starting point to one
ordinary point is exactly one fiber $\rho^{-1}(g_{p_0\to p})$ of the operator
evaluation map;
under the charge projection it fans out into a horizontal segment of the
\ACS{} at fixed multiplicative charge, while under the Reeb quotient it
terminates at distinct heights of one Reeb fiber, the height being the consumed
additive charge.  Second, the rescaled contact form factorizes as
$e^{-M}\alpha=\mathrm d(e^{-M}a)-e^{-M}\mathrm dA$, so the charge projection and
the Reeb quotient are Legendre-conjugate readings of one contact curvature: the
order defect appears as a tearing $\Delta a$ of the assignment fiber in the
first and as a Reeb drift $\Delta A=-\mu\lambda\,\mathrm{Area}_g(\Omega)$ in the
second.  Third, for a pair of bivaluations, uniform motion (constant kernel
coordinate $b$) telescopes to holonomy $1$ around every closed loop, whereas an
acceleration that moves $b$ produces the four-step cross-ratio holonomy
$\Lambda(\square)=(1+a_2b_1)(1+a_1b_2)\big/\bigl((1+a_1b_1)(1+a_2b_2)\bigr)\neq1$,
with an exact finite-field witness $\Lambda=4$ over $\mathbb F_7$.  Fourth, the
incidence-coupled transport $\mathcal T(e)\in\Aff(1,K)$ reduces on the Borel
sector to the affine cocycle of Paper~0 and retains a non-telescoping
multiplicative component off it.  Fifth, the Coulomb radial action is the Reeb
holonomy of a closed horizontal loop, and the monodromy condition on the
quantized Reeb circle returns the Bohr spectrum
$E_n=-m\alpha_e^2/(2\hbar^2n^2)$.

Five further steps require a declared dictionary, and are stated as hypotheses
rather than theorems: the identification of the charged additive step
$\mathsf X_{-e}$ with a physical charge $-e$ and of the reflection
$a\mapsto-a$ with charge conjugation; the identification of the depth exponent
$e^{-M}$ with the Coulomb kernel $\ell_0/r$; the identification of the contact
pairing $\omega(D_u,D_v)=\mu\lambda$ with a bilinear interaction energy
$\kappa=q_1q_2$; the identification of a velocity with the tilt of a history; and the
identification of a nonvanishing projective holonomy with radiated power.  Under
those five hypotheses the note recovers the Coulomb
sign rule (like charges repel, unlike attract), the Amp\`ere coefficient
$(1-\mathbf v_1\!\cdot\!\mathbf v_2/c^2)$ of two uniformly moving charges, a
no-radiation statement for uniform motion with a radiative obstruction for
accelerated motion, the $B_\infty$ versus $B_\infty J B_\infty$
separation of M\o ller and Bhabha scattering, and the hydrogenic spectrum.  No
statement of quantum electrodynamics is proved here, and none is claimed.
\end{abstract}

\noindent\textbf{Keywords:}
arithmetic expression geometry; contact structure; Legendre duality;
torsion--Stokes formula; accumulative commutative space; projective holonomy;
bivaluation; idempotent transport; Coulomb law; Amp\`ere force; Bohr--Sommerfeld
quantization

\bigskip

% ---------------------------------------------------------------------
\begin{cnread}[title={\textbf{中文导读（全篇）}}]
本稿把 \texttt{electrodynamics/} 下五份讨论（\texttt{费曼图.md}、
\texttt{勒让德形式.md}、\texttt{电动力学讨论-01/02/03.md}）整理成一篇正式的
英文文稿，保留原有的物理叙事与论证次序，只做三件事：去掉对话性文字、
统一记号（与 \AEG{} 手稿的符号一致）、并把每一条断言分清来源——是
Paper~0/Paper~IV 已证的定理、是显式声明的「物理字典假设」、还是本稿重新
算过并给出可复核数值的结论。

全篇的主线是一条：\textbf{从理想点 $q$ 到普通点 $p$ 的穿线簇}。
第二、三节把它翻译成算子历史纤维，并在 \ACS{} 与切触空间两个投影下读出
「裸荷发散」与「Reeb 纤维作用量残差」；第四节由此得到库仑定律的符号规律；
第五、六节引入方向角 $\theta$（匀速）与极点跳动（加速），得到安培系数
$1-\mathbf v_1\!\cdot\!\mathbf v_2/c^2$ 与「匀速不辐射、加速才辐射」的
代数判据；第七节把该判据写成完整的射影挠率传递联络 $\mathcal T(e)$，
并在 $\mathbb F_7$ 上给出精确见证 $\Lambda=4$；第八、九节给出
M\o ller/Bhabha 的 Bruhat 扇区对应与玻尔–索末菲能谱 $E_n$。

需要读者留意的是\textbf{分界线}：第一到第五节里真正属于 \AEG{} 内部、
不依赖任何物理解释的数学（算子纤维、勒让德对偶、双重撕裂、交叉比和乐、
Borel 退化、$\mathbb F_7$ 见证、径向作用量）都是可证的；而「荷步 = 电荷」
「$e^{-M}=\ell_0/r$」「$\mu\lambda$ 配对 = 相互作用能」「速度 = 世界线倾角」
「和乐不为 $1$ = 有辐射」这五条是\textbf{字典假设}，第十节与各节末尾逐条列出。它们支撑的
物理结论（库仑、安培、辐射、束缚态）因此都是\textbf{条件性}的：
在字典成立时成立，而不是从 \AEG{} 无条件推出。全文所有可算的恒等式都在
附录~\ref{app:register} 里列了复核脚本，29 项全部通过。
\end{cnread}

\newpage
